\documentclass[11pt,a4paper]{article}
\usepackage[T1]{fontenc}
\usepackage[margin=1in]{geometry}
\usepackage{amsmath,amssymb,mathtools,amsthm}
\usepackage{booktabs}
\usepackage{graphicx}
\usepackage{xcolor}
\usepackage[colorlinks=true,linkcolor=blue!60!black,citecolor=blue!60!black,urlcolor=blue!60!black]{hyperref}
\usepackage{microtype}
\usepackage{enumitem}
\usepackage[most]{tcolorbox}
\usepackage{listings}
\lstset{basicstyle=\ttfamily\scriptsize,breaklines=true,columns=fullflexible,keepspaces=true,frame=single,framerule=0.3pt,xleftmargin=4pt}
% generated by exploration/exciton/python/phase1_paper_numbers.py -- do not edit
\newcommand{\SzeroAdamsRhs}{400}
\newcommand{\SzeroAdamsErr}{$9.76\times10^{-10}$}
\newcommand{\SzeroBdfErr}{$3.11\times10^{-7}$}
\newcommand{\KAaRel}{$1.35\times10^{-8}$}
\newcommand{\KAbRel}{$1.08\times10^{-15}$}
\newcommand{\KAcRel}{$7.23\times10^{-15}$}
\newcommand{\KAtwoRel}{$8.06\times10^{-13}$}
\newcommand{\KAfourRel}{$8.75\times10^{-14}$}
\newcommand{\KAthreeHess}{$4.97\times10^{-9}$}
\newcommand{\KAcAtwoRel}{$2.25\times10^{-13}$}
\newcommand{\KAthreeE}{$5.23\times10^{-14}$}
\newcommand{\KAthreeF}{$6.39\times10^{-16}$}
\newcommand{\FFworstRel}{$1.68\times10^{-12}$}
\newcommand{\FFBcmT}{55.9}
\newcommand{\FFBcRel}{$1.57\times10^{-12}$}
\newcommand{\FFspinNum}{0.730}
\newcommand{\FFspinCF}{0.7254}
\newcommand{\FFstarts}{600}
\newcommand{\SandA}{44.7}
\newcommand{\SandB}{14.2}
\newcommand{\SandC}{4.64}
\newcommand{\SandD}{2.87}
\newcommand{\SandE}{1.87}
\newcommand{\SandF}{1.44}
\newcommand{\SandNA}{22.4}
\newcommand{\SandNB}{13.0}
\newcommand{\SandNC}{10.1}
\newcommand{\SandND}{8.26}
\newcommand{\PtwoAdev}{$1.85\times10^{-14}$}
\newcommand{\PtwoAom}{$1.57\times10^{-15}$}
\newcommand{\PtwoBom}{$5.77\times10^{-11}$}
\newcommand{\PtwoBre}{$2.39\times10^{-11}$}
\newcommand{\PtwoCom}{$6.33\times10^{-9}$}
\newcommand{\PtwoCn}{16}
\newcommand{\LEANNEW}{Lean 4.34.1 with Mathlib tag v4.34.1 (commit d13f23b7)}
\newcommand{\NRuns}{2860}
\newcommand{\NRunsCM}{2060}
\newcommand{\NRunsN}{800}
\newcommand{\NCases}{21}
\newcommand{\NNotConv}{1404}
\newcommand{\NLtwoViol}{0}
\newcommand{\NCommTotal}{9}
\newcommand{\NCommAttained}{8}
\newcommand{\NMinRatio}{$-1.1\times10^{-13}$}
\newcommand{\NFrusMin}{$7.9\times10^{-7}$}
\newcommand{\NFrusMax}{$8.6\times10^{-3}$}
\newcommand{\NNoneBest}{0.663}
\newcommand{\NNoneWorst}{0.724}
\providecommand{\NPARTABSTRACT}{\PENDING{N-particle summary}}
\providecommand{\VERIFYSTATUS}{}

\newcommand{\PENDING}[1]{\textcolor{red!70!black}{[pending: #1]}}
\newtcolorbox{keybox}[1][]{colback=blue!3,colframe=blue!45!black,boxrule=0.5pt,left=4pt,right=4pt,top=3pt,bottom=3pt,#1}
\newtcolorbox{failbox}[1][]{colback=red!3,colframe=red!45!black,boxrule=0.5pt,left=4pt,right=4pt,top=3pt,bottom=3pt,#1}
\newtheorem{lemma}{Lemma}
\newtheorem{proposition}{Proposition}
\newtheorem{theorem}{Theorem}
\newcommand{\elat}{e_{\mathrm{lat}}}
\newcommand{\eH}{e_{\mathrm{H}}}
\DeclareMathOperator{\Adm}{Adm}

\title{\textbf{An optimality theorem and an exciton-fluid interaction:}\\
kernel-checked lemmas, a four-flavour model theorem, and pre-registered N-particle tests on rusty-SUNDIALS}
\author{Xavier Callens\\ \small SocrateAI Research\\
\small \url{https://github.com/xaviercallens/SocrateAI-Scientific-QuantumFluids}}
\date{October 2026}

\begin{document}
\maketitle

\begin{abstract}\noindent
A planar universal-optimality theorem for completely monotone pair interactions (OpenAI, September 2026: AI-generated, not yet
refereed, Lean-formalised and accepted by the Comparator tool) says that the triangular lattice minimises the energy per particle among
all locally finite planar configurations of centred density one. We ask what it says about the interaction between the dipolar interlayer
excitons of a 2026 equilibrium exciton-condensate experiment. We prove in Lean~4, with standard axioms only, that differencing preserves
complete monotonicity, so that the direct exciton--exciton kernel built from \emph{any} admissible single-layer potential is covered;
that the Hartree coefficient is a telescoping integral (the capacitor formula); and that the uniform density minimises a
positive-semidefinite translation-invariant interaction. The consequence is a classical lower bound on the energy per particle which,
at the experimental densities, lies $8$--$22$ times below the Hartree energy: it is structural, not predictive. We prove six
statements about the four-flavour mean-field model of the experiment and reproduce them to \FFworstRel{} with CVODE flows.
We then test the theorem numerically, in a pre-registered design, with gradient flows of $N=36$ and $64$ particles in tori
commensurate and incommensurate with the lattice, integrated by rusty-SUNDIALS CVODE, with non-completely-monotone and planted-bug
controls. \NPARTABSTRACT{} We report what went wrong: a lattice-enumeration slip shared by two implementations and caught only by a
closed-form constant, a compute-budget amendment, and a registered stopping rule too short for soft modes. Producer and verifier of
the Lean files are instances of the same model; nothing here has been refereed.
\end{abstract}

\section{Introduction}

In 2026 Qi \emph{et al.}\ reported evidence for two-component Bose--Einstein condensates of dipolar interlayer excitons in a
MoSe$_2$/hBN/WSe$_2$ electron--hole bilayer with an interlayer distance $d\approx 2$~nm~\cite{Qi2026}. The excitons are injected
electrically and are in equilibrium; there are four spin--valley flavours; the evidence is the spin--valley susceptibility, read through a
four-variable mean-field model; the Berezinskii--Kosterlitz--Thouless temperature reaches about $1.8$~K near the exciton Mott density
$n\approx 0.75\times10^{12}\,\mathrm{cm}^{-2}$. The authors state that counterflow transport and interferometry are not available.

In September 2026 OpenAI released a proof, with a Lean formalisation, that the triangular lattice minimises the energy per particle among
\emph{all} locally finite planar configurations of centred density one for \emph{every} smooth completely monotone function of the squared
distance~\cite{OpenAI2026}. This is the planar case of the Cohn--Kumar universal-optimality conjecture~\cite{CohnKumar2007}, whose
dimensions $8$ and $24$ were settled in 2022~\cite{CKMRV2022}.

Dipolar excitons repel each other; the interaction is a function of the distance. The question of this note is therefore a
\emph{dependency check}: which hypotheses of the theorem does the exciton physics meet, which consequences follow, and how large are
they? Our answers are given in Sections~\ref{sec:kernels} to~\ref{sec:fourflavour}; their kernel-checked parts are listed in
Section~\ref{sec:lean}, and numerical tests with the rusty-SUNDIALS solver~\cite{Hindmarsh2005} are in Section~\ref{sec:numerics}.

\begin{keybox}[title={What this note claims and what it does not}]
\textbf{Claims.} (i) A lemma, proved in Lean, that makes the theorem applicable to the exciton--exciton direct kernel; (ii) the size of the
resulting lower bound relative to the Hartree upper bound at the experimental densities; (iii) six statements about the experiment's own
model, proved in Lean and reproduced numerically; (iv) pre-registered numerical tests of the theorem on small tori, with controls.\\
\textbf{Does not claim.} The theorem is AI-generated and unrefereed; we use it as a hypothesis and test it only on finite systems, which
can fail to falsify it but cannot prove it. The Lean files were produced and checked by instances of the same model (producer
$\neq$ verifier, not human review). Nothing here describes the experiment beyond the authors' own model, whose parameters
$g_X$ and $\Delta$ are phenomenological; agreement of magnitude with the data is not independent evidence. The bound is not tight in the experimental
window. No claim is made about the solver being the best of its kind. Neither the experimental group nor the authors of the
theorem have reviewed this note.
\end{keybox}

\section{The theorem and its conventions}\label{sec:theorem}

Let $g\colon(0,\infty)\to[0,\infty)$ be smooth with $(-1)^r g^{(r)}(t)\ge0$ for all $r\ge0$, $t>0$ (``admissible'', $g\in\Adm$). For a
locally finite $C\subset\mathbb R^2$ write $N_R=\#(C\cap B_R)$, with $B_R$ the closed disk about the origin, and
\begin{equation}
E_g(C)=\liminf_{R\to\infty}\frac1{N_R}\sum_{\substack{x,y\in C\cap B_R\\ x\ne y}} g(|x-y|^2),
\end{equation}
the sum running over \emph{ordered} pairs. If $N_R/(\pi R^2)\to1$ (centred density one), the theorem states
$E_g(C)\ge\sum_{a\in A\setminus 0}g(|a|^2)=E_g(A)$, with $A$ the triangular lattice of covolume one, in $[0,\infty]$.
Four features matter for every use.
\begin{enumerate}[nosep]
\item \emph{Ordered pairs.} The physical energy per particle is half of $E_g$: $\elat(\rho)=\tfrac12\sum_{a\in A_\rho\setminus0}V(|a|)$,
with $A_\rho=\rho^{-1/2}A$. A factor-two slip is the likeliest error downstream; Section~\ref{sec:numerics} plants one on purpose.
\item \emph{Centred disks and a $\liminf$.} The statement is about infinite configurations. A bound for $N$ particles in a
torus needs a periodic-extension argument (Proposition~\ref{prop:sandwich}), elementary on paper and not formalised here.
\item \emph{Value, not minimisers.} The minimiser is not unique in this formulation: a single vacancy changes neither the density nor the
$\liminf$. Uniqueness of the triangular lattice is not part of the theorem.
\item \emph{Class, not each potential.} The proof constructs sharp Fourier minorants for every Gaussian and integrates; a sharp auxiliary
function for each completely monotone potential is not claimed~\cite{OpenAI2026}.
\end{enumerate}
The Lean statement is \texttt{OAI.AtomicTriangular.universal\_energy\_minimum}; the Comparator tool accepted its proof with
axioms $\{\text{propext},\text{Quot.sound},\text{Classical.choice}\}$ in a 140-minute run on another machine (nanoda not run; statement
adequacy not audited by a human; a statement-fidelity audit against the definitions of~\cite{CKMRV2022} was done by an AI
instance)~\cite{LeanMaster2026}. Predecessors are Montgomery's theorem among lattices~\cite{Montgomery1988} and the dimension-$8$ and
$24$ results~\cite{CKMRV2022}.

\section{Kernels, a lemma, and a lower bound}\label{sec:kernels}

\paragraph{The lemma.} Two interlayer excitons, each with its electron in one layer and its hole in the other at distance $d$,
interact directly through $V_{xx}(r)=2[V(r)-V(\sqrt{r^2+d^2})]$ in terms of the single-layer electron--electron potential $V$
(the electron--hole term is the second). In the variable $t=r^2$ this is $2[g(t)-g(t+d^2)]$.

\begin{lemma}[differencing]\label{lem:diff}
If $g\in\Adm$ and $c>0$ then $t\mapsto g(t)-g(t+c)\in\Adm$.
\end{lemma}
\begin{proof}
Put $h_r=(-1)^rg^{(r)}\ge0$. Then $h_r'=-h_{r+1}\le0$, so each $h_r$ is non-increasing. Hence
$(-1)^r\big(g-g(\cdot+c)\big)^{(r)}(t)=h_r(t)-h_r(t+c)\ge0$ for all $r\ge0$ (for $r=0$ this is non-negativity).
\end{proof}

No Bernstein-type theorem is needed. With $g(t)=t^{-1/2}$ (the Coulomb case of the Riesz family, admissible by an explicit formula for the
derivatives) the lemma gives the bilayer dipole kernel $2[t^{-1/2}-(t+d^2)^{-1/2}]$, which decays as $d^2t^{-3/2}$ and has a finite lattice
energy. The same lemma covers every exciton--exciton kernel built from an admissible single-layer potential, in particular the
Keldysh--Rytova potential, which is a positive mixture of exponentials $e^{-ur/r_0}/\sqrt{1+u^2}$ by the integral representation
of $H_0-Y_0$~\cite{DLMF} (verified numerically here to $3\times10^{-31}$, not proved in Lean) and each exponential is completely monotone in
$r^2$; and a Yukawa potential. The potential of a dual-gated device is a sum of image charges with alternating signs; its complete monotonicity
is \emph{not} established here: we checked the sign pattern numerically up to order $5$ for gate distances $5$, $7.5$ and $10$~nm,
which is evidence and not a proof.

\paragraph{Kernel-checked.} In Lean~4 (Section~\ref{sec:lean}): Lemma~\ref{lem:diff}; that the bilayer kernel is admissible given the
Riesz case $s=1$ (which is a corollary proved by the LeanMaster contribution~\cite{LeanMaster2026}); and the negative control that
$e^{-t^2}$ (the cluster-crystal kernel of the generalised exponential model) is \emph{not} admissible.

\paragraph{The Hartree coefficient.} For admissible $g$ with $g\to0$ at infinity,
\begin{equation}\label{eq:hartree}
\int_{\mathbb R^2}\!\big[g(|x|^2)-g(|x|^2+d^2)\big]\,d^2x=\pi\!\int_0^{d^2}\!g(t)\,dt
\end{equation}
(telescoping; Lean: \texttt{telescoping\_integral}). For $g=2t^{-1/2}$ this is $4\pi d$, in units $e^2/4\pi\varepsilon_0\varepsilon=1$:
the capacitor formula $\tilde U(0)=e^2d/\varepsilon_0\varepsilon$, and the authors' $g_H=8\pi d$ in Rydberg units
(Lean: \texttt{bilayer\_hartree}).

\paragraph{Mean field.} A completely monotone function of $r^2$ is positive definite in every dimension~\cite{Schoenberg1938}, so its Fourier
transform is non-negative. For a positive-semidefinite translation-invariant interaction on a finite abelian group, the uniform density
minimises $\sum_{x,y}n(x)n(y)U(x-y)$ at fixed mass (Lean: \texttt{uniform\_minimises}, with positive semidefiniteness as a hypothesis).
Consequently a Gross--Pitaevskii calculation with such a kernel cannot crystallise: the uniform state is its ground state, and the
Bogoliubov spectrum has no roton. That is a statement about mean-field theory, not about the quantum fluid.

\begin{proposition}[sandwich, on paper]\label{prop:sandwich}
For $N$ bosons on a torus with the periodised completely monotone kernel $V_L$, the ground-state energy per particle satisfies
$\elat(\rho)\le e_0(\rho)\le\eH(\rho)(1-1/N)$, with $\eH=\tfrac12\rho\tilde U(0)$.
\end{proposition}
\begin{proof}[Sketch]
The periodic extension of any $N$-point configuration has centred density $\rho$ and $E_g$ equal to the per-cell average; the theorem
and the non-negativity of the kinetic energy give the lower bound; the constant orbital gives the upper bound. The periodic-extension
step is elementary and is not formalised here.
\end{proof}

\begin{table}[ht]\centering
\begin{tabular}{@{}lrrrrrr@{}}
\toprule
$\rho d^2$ & $10^{-3}$ & $10^{-2}$ & $0.1$ & $0.3$ & $1$ & $3$\\
$\eH/\elat$ & \SandA & \SandB & \SandC & \SandD & \SandE & \SandF\\
\bottomrule
\end{tabular}\\[4pt]
\begin{tabular}{@{}lrrrr@{}}
\toprule
$n$ ($10^{12}$~cm$^{-2}$), $d=2$~nm & $0.1$ & $0.3$ & $0.5$ & $0.75$\\
$\eH/\elat$ & \SandNA & \SandNB & \SandNC & \SandND\\
\bottomrule
\end{tabular}
\caption{Hartree energy over classical lattice energy for the undamped bilayer kernel (independent numpy and Rust sums agree to
$9\times10^{-14}$). In the experimental window the lower bound sits one decade below the Hartree energy: the fluid is deep in the quantum
regime, the bound is tight only in the crystal limit, which the Mott transition pre-empts.}
\end{table}

\section{The experiment's own model}\label{sec:fourflavour}

With flavours $1,2=\mathrm{KK},\mathrm{K'K'}$ (intravalley) and $3,4=\mathrm{KK'},\mathrm{K'K}$ (intervalley), the authors' energy is
\begin{equation}
H(n)=\sum_i(E_i-\mu)n_i+\tfrac12(g_H+g_X)\Big(\sum_in_i\Big)^2-g_X(n_1n_2+n_3n_4),\qquad n_i\ge0,
\end{equation}
with $E_{1,2}=\pm(g_v-g_c)\mu_BB-\Delta$, $E_{3,4}=\mp(g_c+g_v)\mu_BB$, $g_H=8\pi d$, $g_c\approx3$, $g_v\approx6$ and the
phenomenological $g_X=1$, $\Delta=1~\mu$eV (Rydberg and Bohr-radius units; arXiv version 1). A four-variable quadratic energy is exactly what a
proof assistant handles. We proved, in Lean, for $g_X>0$:
\begin{enumerate}[nosep,label=\textbf{T\arabic*.}]
\item a minimiser condenses at most one exchange pair, $\{1,2\}$ or $\{3,4\}$, except on the degenerate set $E_i=E_j$;
\item on $\{1,2\}$: $n_2-n_1=2(g_v-g_c)\mu_BB/g_X$ and $N=2(\mu+\Delta)/(2g_H+g_X)$;
\item on $\{3,4\}$: $n_3-n_4=2(g_c+g_v)\mu_BB/g_X$ and $N=2\mu/(2g_H+g_X)$;
\item the grand potentials $\Omega_A=-(\mu+\Delta)^2/(2g_H+g_X)-\beta_A^2/g_X$ and
$\Omega_B=-\mu^2/(2g_H+g_X)-\beta_B^2/g_X$, $\beta_A=(g_v-g_c)\mu_BB$, $\beta_B=(g_c+g_v)\mu_BB$;
\item the first-order transition between them at $(\mu_BB_c)^2=g_X\Delta(2\mu+\Delta)/\big(4g_cg_v(2g_H+g_X)\big)$;
\item for $g_X<0$ every minimiser is single-component (the Hartree--Fock ferromagnet; the negative control).
\end{enumerate}
T2 and T3 are the authors' printed formulas (a confirmation); T1 and T4--T6 are not printed in the main text of the arXiv version. With the
quoted parameters at $n_x=0.5\times10^{12}\,\mathrm{cm}^{-2}$ the closed form gives $B_c=\FFBcmT$~mT; the authors' windows are
$|B|<30$~mT (II$_A$) and $30$--$70$~mT (II$_B$) at that density. The agreement of magnitude is expected, since $g_X$ and $\Delta$ were chosen
to describe the data. The statement $B_c\propto\sqrt n$ at fixed $\Delta$ reproduces the trend they report.

\section{Lean formalisation}\label{sec:lean}

\begin{table}[ht]\centering\small
\begin{tabular}{@{}llp{7.2cm}@{}}
\toprule
file & role & named results\\
\midrule
\texttt{ExcitonX1.lean} & Lemma~\ref{lem:diff} & \texttt{admissible\_sub\_shift}, \texttt{admissible\_const\_mul}, \texttt{bilayerDipole\_admissible} (given $s=1$), \texttt{gem4\_not\_admissible}\\
\texttt{FourFlavour.lean} & T1--T6 & \texttt{support\_in\_one\_pair}, \texttt{IIA\_polarisation}, \texttt{IIB\_polarisation}, \texttt{intravalley\_density}, \texttt{grand\_potential\_gap}, \texttt{critical\_field}, \texttt{single\_component\_of\_neg\_gX}\\
\texttt{MeanField.lean} & Eq.~\eqref{eq:hartree}, mean field & \texttt{uniform\_minimises}, \texttt{telescoping\_integral}, \texttt{bilayer\_hartree}, \texttt{variational\_lower\_bound}\\
\texttt{GradientFlow.lean} & energy monitor & \texttt{energy\_antitone}, \texttt{energy\_monotone\_ascent} (control)\\
\texttt{FourFlavourNegativeControl.lean} & control & wrong-sign polarisation, \emph{fails to compile by design}\\
\bottomrule
\end{tabular}
\caption{The Lean files (\texttt{exploration/exciton/lean/}). Every named result depends only on
\{propext, Classical.choice, Quot.sound\}; no \texttt{sorry}. Compiled under Lean 4.34.0-rc2 / Mathlib 85e3a25 (the pin of the
QuantumFluids library) and under \LEANNEW{} (Lean 4.34.1, Mathlib tag v4.34.1 = d13f23b7, the pin of the upstream project).}
\end{table}
\VERIFYSTATUS{}

\section{Numerical tests on rusty-SUNDIALS}\label{sec:numerics}

All experiments were \emph{pre-registered} before they ran (\texttt{docs/designs/EXCITON\_FLUID\_PHASE1\_PREREG.md}, committed with
independent reference numbers produced by a separate numpy/mpmath script before the Rust code existed); gates are fixed there and two
amendments, with their reasons and dates, are recorded in the same file and below. Results are read as registered, failures included.

\subsection{Instrument}

The code is a small Rust crate on rusty-SUNDIALS tree \texttt{996aaf0706e1} (release v11.6.0, commit \texttt{5db8041}), using the
\texttt{cvode} crate with the BDF method (orders 1--5), relative tolerance $10^{-10}$ and an analytic Jacobian. Gate S0: the crate's own
tests pass ($22$ of $22$, including the Adams-order and tight-tolerance tests); on $y'=-y$ to $t=10$ the Adams method at $\mathrm{rtol}=10^{-10}$
needs \SzeroAdamsRhs{} right-hand-side evaluations for a maximum relative error \SzeroAdamsErr{} (the pre-fix first-order build of the
shared Python environment needed about $2\times10^{6}$), and BDF at $\mathrm{rtol}=10^{-8}$ reaches \SzeroBdfErr.

\subsection{Known answers, and the first failure}

The energy per particle of $N$ points in a rectangular torus, with periodic images summed and the images of a point with itself included
(as in the periodic extension the theorem is applied to), is $E=\frac1{2N}\sum_{(i,j,m)\ne(i,i,0)}g(|x_i-x_j+L\!\circ\! m|^2)$; the
lattice energy is $\elat(\rho)=\tfrac12\sum_{a\ne0}g(|a|^2)$. The gates were: (KA-1a) $\sum' r^{-3}$ over the unit-spacing triangular
lattice against $6\zeta(3/2)L(3/2,\chi_{-3})=11.0341757349\ldots$; (KA-1b) the Gaussian lattice energy against the Jacobi-theta closed form;
(KA-1c) the lattice energies of all exactly summable kernels against an independent numpy sum; (KA-2) the Hartree integral
$\int V\,d^2r=4\pi d$ by quadrature; (KA-3) the Rust energy and force against an independent numpy implementation at $35$ random
configurations and the analytic Hessian against a finite difference of the force; (KA-4) the table of $\eH/\elat$.

\begin{failbox}[title={KA-1a failed as registered, first}]
The relative errors were $1.8\times10^{-4}$, $1.0\times10^{-4}$, $5.4\times10^{-5}$ at cutoffs $R=200,400,800$, halving with $R$. The lattice
enumeration of both the numpy reference and the Rust code looped over $|m|,|k|\le R/a+3$ and missed the caps $|y|>\tfrac{\sqrt3}2R$ of
the disk. Two implementations by the same author shared the slip; only the closed-form constant exposed it. Both were corrected, the
references regenerated and the gates re-run unchanged (amendment A1). The exactly summable references changed by less than $2\times10^{-16}$,
the sandwich table by at most $3\times10^{-4}$.
\end{failbox}

After the correction the gates read: KA-1a \KAaRel{} at $R=800$ (gate $10^{-6}$); KA-1b \KAbRel{}; KA-1c \KAcRel{} (and \KAcAtwoRel{} at the
reduced cutoffs of amendment A2); KA-2 \KAtwoRel{}; KA-4 \KAfourRel{}; KA-3 energy \KAthreeE{}, force \KAthreeF{} (max-norm, relative),
Hessian against finite differences \KAthreeHess{}; a Hessian with the sign flipped (planted) is detected.

\subsection{The four-flavour flow}

The imaginary-time flow $\dot\psi_i=-\psi_i\,\partial H/\partial n_i$, $n_i=\psi_i^2$, was integrated with CVODE (BDF, $\mathrm{rtol}=10^{-12}$)
from \FFstarts{} random positive starts at $B=0.01$, $0.04$ and $0.2$~T. Every converged final state has its support inside one exchange
pair (T1: no cross-pair supports among $\FFstarts{}$); the polarisations and total densities of T2, T3 are reproduced to
\FFworstRel{} (gate $10^{-9}$); the field at which the grand potentials of the converged II$_A$ and II$_B$ states are equal is
$B_c=\FFBcmT$~mT, equal to the closed form to \FFBcRel{} (gate $10^{-8}$); and starting from II$_A$ with a seed in flavour~2 the state
leaves the intravalley pair at \FFspinNum{}~T against the closed-form spinodal \FFspinCF{}~T (gate $5\,\%$). The large metastability window
is a property of the mean-field model, not a prediction of the experiment. The grand potential decreases along every flow.

\subsection{N-particle gradient flows}
\PENDING{results of the N-particle experiments}

\section{Limitations and what we do not know}
\PENDING{limitations}

\section*{Reproducibility}
\PENDING{reproducibility}

\bibliographystyle{unsrt}
\bibliography{refs_exciton}
\end{document}
